\section{A Survey of Graph Learning for Wireless Networks, by Operator}
\label{sec:survey}

Having fixed the propagation-operator view, we use it to organize the landscape.
Fig.~\ref{fig:taxonomy} is the map: every method we discuss is a choice of operator
$\Prop(\Lap)$, optionally wrapped to handle a time-varying graph. This section surveys
the four branches and their use in wireless and non-terrestrial networks; the depth on
each is intentionally uneven, favoring the operators we later build and test.

\begin{figure*}[t]
\centering
\resizebox{\textwidth}{!}{%
\begin{tikzpicture}[
  every node/.style={font=\footnotesize},
  root/.style={box, fill=navy!12, draw=navy, font=\footnotesize\bfseries},
  fam/.style={hot, align=center},
  leaf/.style={plain, align=center, font=\scriptsize},
  edge from parent/.style={draw=steel, thick, -{Stealth[length=2mm]}},
  level 1/.style={sibling distance=64mm, level distance=14mm},
  level 2/.style={sibling distance=22mm, level distance=16mm},
]
\node[root] {Graph learning as a\\propagation operator $\Prop(\Lap)$}
  child {node[fam] {Classical\\(local / diffusive)}
    child {node[leaf] {spatial MP\\GCN, GAT,\\GraphSAGE, GIN}}
    child {node[leaf] {spectral\\ChebNet,\\graph wavelets}}
    child {node[leaf] {diffusion\\GDC, APPNP\\$e^{-t\Lap}$}}
  }
  child {node[fam] {Quantum\\(ballistic / unitary)}
    child {node[leaf] {quantum walks\\CTQW $e^{-it\Lap}$,\\QWNN}}
    child {node[leaf] {QGNN / PQC\\equivariant QGC,\\quantum kernels}}
  }
  child {node[fam] {Quantum-inspired\\(runnable today)}
    child {node[leaf] {$|e^{-it\Lap}|^2$ as a\\classical layer\\(this tutorial)}}
  }
  child {node[fam] {Dynamic-graph\\wrappers}
    child {node[leaf] {TGN, EvolveGCN,\\DySAT (any\\operator above)}}
  };
\end{tikzpicture}}
\caption{A taxonomy of graph learning organized by the propagation operator
$\Prop(\Lap)$. Classical operators are local (spatial message passing) or diffusive
(spectral / heat); quantum operators are unitary and ballistic (continuous-time quantum
walks, parameterized quantum graph circuits); the quantum-inspired branch runs the
quantum-walk kernel on classical hardware; dynamic-graph wrappers compose with any
operator to handle a time-varying topology. This tutorial threads the classical,
quantum, and quantum-inspired branches onto one axis.}
\label{fig:taxonomy}
\end{figure*}


\subsection{Classical operators}
The spatial branch of message passing, GCN~\cite{kipf2017gcn},
GraphSAGE~\cite{hamilton2017graphsage}, GAT~\cite{velickovic2018gat},
GIN~\cite{xu2019gin}, and the general message-passing framework~\cite{gilmer2017mpnn},
shares a $K$-hop horizon and is surveyed comprehensively
elsewhere~\cite{wu2021gnnsurvey,zhou2020gnnreview,bronstein2017geometric}. The spectral
branch, rooted in graph signal processing~\cite{shuman2013gsp,ortega2018gsp} and spectral
graph theory~\cite{chung1997spectral}, expresses a layer as a function of $\Lap$:
ChebNet~\cite{defferrard2016chebnet} and graph wavelets~\cite{hammond2011wavelets} are
polynomials of the spectrum, while the diffusion branch
(GDC~\cite{gasteiger2019diffusion}, APPNP~\cite{gasteiger2019appnp},
GCNII~\cite{chen2020gcnii}) uses the heat kernel $e^{-t\Lap}$ or its rational cousins to
propagate globally. The recurring failure mode of the local branch, the inability to
carry information past the depth horizon, is studied as over-smoothing~\cite{li2018deeper}
and over-squashing~\cite{alon2021bottleneck,topping2022oversquashing}; it is exactly the
deficiency our case study exposes.

\subsection{Quantum and quantum-inspired operators}
The quantum branch replaces dissipative diffusion by unitary evolution. Continuous-time
quantum walks~\cite{farhi1998quantumwalk,childs2004quantumwalk}, reviewed
in~\cite{kempe2003qwalk,venegas2012quantumwalks,portugal2013book}, underlie quantum
speedups~\cite{childs2003exponential} and are computationally
universal~\cite{childs2009universal}; their learning incarnation is the quantum-walk
neural network~\cite{dernbach2019qwnn}. The parameterized-circuit branch encodes graph
features into qubits and learns with variational
circuits~\cite{mitarai2018qcl,farhi2018classification,cerezo2021vqa}: quantum graph
neural networks~\cite{verdon2019qgnn}, equivariant quantum graph
circuits~\cite{mernyei2022equivariant}, and quantum graph kernels~\cite{henry2021quantumkernel}.
Their appeal is the exponential-dimensional feature
map~\cite{schuld2019featurespace,havlicek2019supervised,schuld2018book} and, in some
settings, provable data or speed advantages~\cite{huang2021power,liu2021rigorous};
their obstacles are NISQ noise~\cite{preskill2018nisq}, barren
plateaus~\cite{mcclean2018barren}, and limited qubits, surveyed
in~\cite{biamonte2017qml,dunjko2018review}. The quantum-inspired branch sidesteps the
hardware by running the quantum-walk kernel $|e^{-it\Lap}|^2$ as a classical layer; it is
the operator we build and evaluate, and Table~\ref{tab:opfamilies} places it against the
others.

\subsection{Handling the time-varying topology}
Any operator above can be wrapped to track a changing graph. Dynamic-graph
methods~\cite{kazemi2020dynamicsurvey} such as TGN~\cite{rossi2020tgn},
EvolveGCN~\cite{pareja2020evolvegcn}, and DySAT~\cite{sankar2020dysat} learn over a
sequence of snapshots; for LEO networks the future graph is moreover known from
ephemeris, a structure most of these methods do not exploit and that we flag as an
opportunity in Section~\ref{sec:roadmap}.

\begin{table}[t]
\centering
\caption{Operator families on the propagation axis. $N$ nodes, $|E|$ edges,
$K$ layers/hops, $d$ feature width.}
\label{tab:opfamilies}
\renewcommand{\arraystretch}{1.25}
\footnotesize
\begin{tabular}{@{}p{0.20\columnwidth}p{0.20\columnwidth}p{0.16\columnwidth}p{0.18\columnwidth}@{}}
\toprule
Family & Operator $\Prop$ & Reach / spread & Cost (layer) \\
\midrule
Spatial MP & $\hat{\Adj}$ (1-hop) & local, $K$ hops & $O(|E|d)$ \\
Diffusion & $e^{-t\Lap}$ & global, $\sim\!\sqrt{t}$ & $O(N^3)$ eig$^\dagger$ \\
Quantum walk & $e^{-it\Lap}$ (unitary) & global, $\sim\! t$ & $O(N^3)$ eig$^\dagger$ \\
Quantum-inspired & $|e^{-it\Lap}|^2$ & global, $\sim\! t$ & $O(N^3)$ eig$^\dagger$ \\
True QGNN/PQC & PQC on qubits & Hilbert-space & NISQ-limited \\
\bottomrule
\end{tabular}
\\[2pt]
{\footnotesize $^\dagger$ amortized over layers; reducible to $O(|E|)$ per
matrix-vector product with Chebyshev/Krylov approximation (Section~\ref{sec:roadmap}).}
\end{table}

\subsection{Graph and quantum learning in wireless and non-terrestrial networks}
On the application side, graph learning is now established for communication
networks~\cite{suarezvarela2023gnn,jiang2022graph}; we group the uses by the task they
place on the graph.

\emph{Routing and traffic engineering.} The earliest and most direct use treats the
network as a graph and learns to predict path-level quantities. RouteNet models per-path
delay, jitter, and loss and uses the learned model to optimize
routing~\cite{rusek2020routenet}; the appeal is that a graph model generalizes across
topologies a per-topology optimizer cannot reuse. This is the regime where the local
versus global operator distinction of this tutorial bites hardest, because routing
potentials are long-range fields on a large-diameter graph.

\emph{Resource management and physical layer.} A second line learns resource-allocation
and power-control policies as functions on the interference or connectivity graph, where
permutation-equivariant graph operators give scalability and transferability that dense
networks lack~\cite{shen2021gnnradio,eisen2020optimal}, frequently combined with
reinforcement learning for sequential control~\cite{mao2016resource}. Here the graph is
often denser and lower-diameter, so a local operator can suffice, a useful contrast to
the routing regime that the operator view makes explicit.

\emph{Non-terrestrial and satellite networks.} LEO
mega-constellations~\cite{delportillo2019technical} are the setting of our case study:
they pose low-latency routing over a fast-moving but ephemeris-predictable
topology~\cite{handley2018delay}, with packet-level simulators available for honest
evaluation~\cite{kassing2020hypatia}. Because the constellation is physically distributed
and bandwidth to the ground is scarce, learning is increasingly pushed on board and
across satellites through federated schemes~\cite{razmi2022board,matthiesen2024federated},
which raises the decentralization question we return to in Section~\ref{sec:roadmap}.

\emph{Quantum methods for networking.} Quantum techniques enter this space along two
routes: quantum search and optimization for discrete network functions such as
virtualization and routing~\cite{botsinis2019quantum}, and the longer-term
quantum-internet vision in which entanglement distribution is itself a network
service~\cite{wehner2018quantum}. Graph-structured quantum learning, the subject of this
tutorial, sits between them: a learned propagation operator rather than a fixed search
subroutine.

Table~\ref{tab:qgraphworks} summarizes representative quantum graph-learning works by how
they encode the graph and what they target. The gap this tutorial fills is the
controlled, honest comparison of these operators on a single dynamic-wireless task, which
the application literature above rarely isolates.

\begin{table}[t]
\centering
\caption{Representative quantum / quantum-inspired graph-learning works by encoding and
target. Sim = classical simulation; HW = quantum hardware.}
\label{tab:qgraphworks}
\renewcommand{\arraystretch}{1.25}
\footnotesize
\begin{tabular}{@{}p{0.27\columnwidth}p{0.30\columnwidth}p{0.18\columnwidth}@{}}
\toprule
Work & Mechanism / encoding & Realization \\
\midrule
CTQW~\cite{childs2004quantumwalk} & unitary $e^{-it\Lap}$, search & theory \\
QWNN~\cite{dernbach2019qwnn} & learned-coin quantum walk & Sim \\
QGNN~\cite{verdon2019qgnn} & Hamiltonian-based PQC & Sim \\
EQGC~\cite{mernyei2022equivariant} & equivariant graph circuit & Sim \\
Quantum kernel~\cite{henry2021quantumkernel} & evolution kernel, qubit array & Sim/HW \\
This tutorial & $|e^{-it\Lap}|^2$ classical layer & Sim (CPU) \\
\bottomrule
\end{tabular}
\end{table}
